\documentclass[12pt]{article}
\usepackage[utf8]{inputenc}
\usepackage{graphicx}
\usepackage{amsmath}
\usepackage{geometry}
\usepackage{setspace}
\usepackage{hyperref}
\usepackage{booktabs}
\usepackage{siunitx}
\usepackage{float}
\geometry{margin=1in}
\setstretch{1.5}

\usepackage[backend=biber,style=ieee]{biblatex}
\addbibresource{references.bib}

\title{Automated Combat Sports Analysis Using Computer Vision}
\author{Jimmy Huynh}
\date{\today}

\begin{document}

\maketitle

\newpage

\begin{abstract}
This thesis presents the design and evaluation of a reference-based system for automated assessment of martial arts techniques. The system leverages pose estimation to extract skeletal keypoints from video recordings and compares user-performed movements against predefined reference sequences. By combining spatial alignment and temporal similarity measures, the system produces a continuous score reflecting the quality of execution and provides structured feedback to the user.

The primary objective of this work is to investigate whether such a system can meaningfully distinguish between different levels of technical proficiency and approximate human expert evaluation. In the current project stage, clip-level trainer outputs were analyzed to assess score behavior across stronger and weaker executions. The codebase also supports the planned next step of comparing trained and amateur users, and comparing system-generated scores with subjective assessments provided by an experienced trainer.

Overall, the findings demonstrate that reference-based pose comparison is a viable foundation for automated martial arts training support, while also identifying key areas for future improvement, including enhanced temporal modeling and incorporation of higher-level performance features.
\end{abstract}

\newpage

\section{Introduction}
Combat sports analysis has traditionally relied on human judgement, which introduces subjectivity and inconsistency. Judges evaluate techniques based on criteria such as precision, impact, balance, and control, but this process is vulnerable to bias and error in fast-moving exchanges \cite{quinn2024automation}. Recent advances in computer vision and deep learning provide an alternative approach by allowing human actions to be analysed directly from video. In combat sports, this creates the possibility of extracting pose information, comparing an athlete's movements against a technical reference, and generating consistent quantitative feedback \cite{wu2024sthka,echeverria2021pose}.

This thesis develops a computer-vision-based system for technique analysis and estimation in combat sports. Rather than treating the problem only as action classification, the implemented system is designed as a pose-driven trainer that evaluates the quality of execution of a target technique. The project combines person tracking, pose estimation, temporal pose sequence processing, multi-angle reference matching, and calibrated pass/fail thresholds. The overall aim is to support training and evaluation by producing both a numerical score and corrective feedback.

A further goal of the project is to support the type of evaluation requested by the supervisor: comparison between amateur and professional users of the program, and comparison between a professional oral score and the program score on a 1--10 scale. The current codebase supports the core scoring logic required for such experiments. This report therefore explains the design of the implemented system, the practical challenges encountered during development, how those challenges were addressed, and how the final experimental results should be structured.

\section{Background and Literature Review}
\subsection{Human Action Recognition and Pose Estimation}
Human Action Recognition (HAR) seeks to identify movements or activities from image or video data. In sports contexts, HAR is especially useful because it can transform raw video into structured information about technique, and performance. However, sports video is difficult to analyse due to occlusion, variable viewpoints, large appearance differences, and rapid motion. These difficulties are even more pronounced in combat sports, where a technique may last only a fraction of a second and body parts are often partially hidden \cite{wu2024sthka}.

Pose estimation addresses part of this problem by reducing video to a set of semantically meaningful keypoints. Instead of relying on raw pixels, a pose-based method represents the athlete as a skeleton, making it easier to compare body structure over time and across performers. Prior work has shown that pose estimation is especially relevant for martial arts and karate because technique quality is strongly linked to joint alignment, body mechanics, and coordinated motion \cite{echeverria2021pose}. This makes pose-based analysis an appropriate basis for a coaching-oriented system.

\subsection{Combat Sports Analytics}
Recent research has explored the use of computer vision for combat sports scoring, athlete tracking, and action recognition. Quinn and Corcoran argue that real-time object detection and pose estimation can be combined to generate meaningful performance statistics in combat sports and potentially reduce the effects of human error in scoring \cite{quinn2024automation}. Similarly, work on optimized YOLO pipelines for combat sports highlights the importance of careful data collection, data augmentation, modular system design, and domain-specific optimization when working with athletic video data \cite{quinn2025yolo}.

At the same time, the literature makes clear that generic video classification alone is not enough. A technically useful system must cope with viewpoint changes, noisy data, short action windows, and the absence of large labelled datasets specific to combat sports \cite{wu2024sthka,quinn2025yolo}. These issues strongly influenced the design decisions in the present project.

\section{System Design}
The implemented project is not a simple end-to-end classifier that takes a video and outputs a label. Instead, it is a modular pose-driven pipeline built around person tracking, pose sequence extraction, reference matching, and calibrated scoring. The central implementation is contained in \texttt{action\_recognition.py}, which uses an Ultralytics YOLOv26 pose model to detect and track people, extract keypoints in COCO-17 format, and maintain temporal histories of those keypoints for each tracked identity \cite{projectActionRecognition}. The code defines support for the labels \textit{fighting stance, jab, cross, hook, uppercut, front kick, roundhouse kick, side kick, back kick, spinning back kick, knee strike, elbow strike,} and \textit{axe kick} \cite{projectActionRecognition}.

The pipeline begins with a video source, which may be a webcam, local file or a YouTube URL. For each frame, the YOLO pose model performs person detection and pose estimation. The extracted poses are associated with persistent track identifiers, allowing the system to follow one athlete over time rather than analysing isolated frames only \cite{projectActionRecognition}. This temporal continuity is important because a technique such as a jab or kick cannot be judged from a single still image.

The implementation also supports optional video classification. The code includes interfaces for pretrained TorchVision video models such as S3D, R3D-18, Swin3D, and MViT, as well as a zero-shot Hugging Face X-CLIP classifier \cite{projectActionRecognition}. However, the project architecture clearly prioritises pose-based scoring. In many of the reference-collection workflows, the video classifier is explicitly disabled so that the system focuses on clean pose extraction and technique comparison instead of generic action labels \cite{projectFrontKickLog,projectSeedLog}.

A central design element is the reference pose library. Canonical techniques are stored as \texttt{.npy} files containing pose sequences. The loader supports both a legacy flat structure and a multi-angle layout of the form \texttt{reference\_poses/<technique>/<angle>.npy} \cite{projectActionRecognition,projectReadme}. This is a significant methodological decision because it recognises that the same technique may look different from different camera angles. Thus this implementation implicitly adds the functionality of adding variations of a technique. Rather than forcing all users to match a single canonical view or movement, the system chooses the best-matching angle and movement from the reference bank \cite{projectActionRecognition}.

The scoring process begins by normalising each pose frame. The body is centred using the hip midpoint when available, and the coordinates are scaled by shoulder distance or a fallback body-based scale \cite{projectActionRecognition}. This reduces the influence of camera distance, translation, and body size. The system then compares the user sequence against the reference sequence using several pose-based measures. These include Dynamic Time Warping for temporal alignment, cosine similarity of flattened pose trajectories, and technique-specific joint angle error using the joints most relevant to punches or kicks \cite{projectActionRecognition}. For striking actions, the code also computes wrist motion energy and wrist return closure, two handcrafted measures intended to capture whether a movement contains both meaningful extension and a realistic return phase \cite{projectActionRecognition}.

The trainer output is therefore richer than a class label. For each evaluated attempt, the system can return a best-matching angle, a numerical score, a score threshold, a correctness decision, and generated feedback \cite{projectActionRecognition}. If the attempt falls below the threshold, the code constructs a ghost-pose overlay from the reference sequence, allowing visual comparison between the observed movement and the expected technique \cite{projectActionRecognition}. This is particularly useful in a coaching scenario because it turns the system into a feedback tool rather than a black-box classifier. Thus this tool will assist a student by identifying unaligned joints and movements, visualized by the correct reference pose, that fits the direction and technique.

The implementation also includes structured storage of experimental artifacts. It can create dedicated run folders containing configuration metadata, metric rows, track summaries, and overlays \cite{projectActionRecognition}. From a research perspective, this is valuable because it improves reproducibility, makes experiments easier to compare, and supports later analysis of failures and edge cases.

\section{Challenges and Solutions}
\subsection{Challenge 1: Building a Reference Dataset}
One of the most difficult parts of the project was not model inference itself, but the creation of a reliable reference library. Unlike benchmark-driven projects that begin with a large labelled dataset, this system depends on curated canonical examples of techniques. This aligns with the literature, which identifies data collection, annotation, and dataset quality as major constraints in combat sports computer vision \cite{quinn2024automation,quinn2025yolo}.

The codebase addresses this challenge through a semi-automated workflow. 

First, \texttt{scrape\_jab\_candidates.py} collects candidate YouTube videos for a target technique using the YouTube Data API and stores them with metadata such as title, duration, and URL \cite{projectScrape}. Second, \texttt{filter\_candidates.py} applies hand-designed keep and reject keyword patterns to filter obvious false positives and non-technique content \cite{projectFilter}. Third, \texttt{generate\_reference\_capture\_commands.py} converts reviewed CSV rows into executable PowerShell commands for reference capture \cite{projectGenerate}. Fourth, \texttt{run\_reference\_collection\_batch.py} executes these capture jobs in batch mode with checks for source diversity, controls on examples per angle, CPU throttling, and cooldown periods \cite{projectRunBatch}. This workflow directly addresses the challenge of collecting high-quality references from messy web video sources.
The most consistent type of data has been extracting self made videos or reviewed youtube videos, as it can be recorded in controlled environments to reduce noise.

\subsection{Challenge 2: Viewpoint Sensitivity}
A technique in 2D pose space can vary substantially with camera placement. A jab observed from the side may appear very different from a jab observed from the front, even when biomechanically correct. Prior literature identifies viewpoint variation as a core difficulty in martial arts and action recognition more broadly \cite{wu2024sthka,echeverria2021pose}.

During development it has been discovered that a side view, the Orthodox (left foot/hand forward, right-handed) fighting stance yields the best results. Technical difficulties can bee see using Soutpaw (right foot/hand forward, left-handed), as this causes hidden joints from specific angles. This causes the normalization to misalign center position and the sizing of the pose figure.

The project responds to this problem with a multi-angle reference design. The project documentation describes a technique-first folder layout in which each technique can store separate angle files such as \texttt{front}, \texttt{left45}, \texttt{right45} \texttt{left}, and \texttt{right} \cite{projectReadme}. During evaluation, the code does not compare an athlete against a single reference template. Instead, it evaluates all available angle references for the target technique and selects the one with the best score \cite{projectActionRecognition}. This substantially improves robustness to viewpoint mismatch.

\subsection{Challenge 3: Temporal Segmentation of Techniques}
A combat technique is not defined by one frame. A punch, for example, includes initiation, extension, peak action, and recovery. If a saved reference window begins too early or ends too late, the resulting sequence contains irrelevant motion and becomes a weak template. This problem is especially important in combat sports because many actions are brief and explosive, but during technical training, often slow movement and prolonged duration.

The project addresses this challenge inside \texttt{action\_recognition.py} through explicit sequence extraction logic. The code supports event-centred extraction around peak wrist extension and a stance-cycle mode based on start, peak, and end thresholds \cite{projectActionRecognition}. It also uses wrist motion energy and wrist return closure as gates before saving a reference, so that only sequences with meaningful motion and plausible return behaviour are accepted \cite{projectActionRecognition}. This is a practical solution to the problem of noisy or poorly segmented reference windows.
Combining this solution with Dynamic Time Warping creates a program, that work for both explosive and slow actions.

\subsection{Challenge 4: Translating Similarity into Correctness}
A raw score is informative, but it does not by itself define what counts as a correct technique. A threshold that is appropriate for one action may be too strict or too lenient for another. This problem becomes more severe when different techniques vary in visibility, temporal stability, and sensitivity to joint noise.

The solution implemented in the project is \texttt{calibrate\_label\_thresholds.py}. This script assumes a validation layout of \texttt{datasets/validation/<technique>/correct/*.npy} and \texttt{datasets/validation/<technique>/incorrect/*.npy}. For each technique, it scores validation sequences against the reference bank, searches across candidate thresholds, and chooses the threshold that maximizes F1 score and then balanced accuracy \cite{projectCalibrate}. It also reports precision, recall, accuracy, balanced accuracy, and the counts of true positives, false positives, true negatives, and false negatives \cite{projectCalibrate}.

\subsection{Challenge 5: Runtime Constraints and Detection Gaps}
The uploaded runtime logs show that pose-based reference collection is feasible on CPU hardware, but not trivial. This project has been run using the following specifications:

\begin{itemize}
    \item Processor	AMD Ryzen 9 5900X 12-Core Processor             (3.70 GHz)
    \item Installed RAM	16,0 GB
    \item System type	64-bit operating system, x64-based processor
    \item NVIDIA GeForce GTX 1660 SUPER
    \item NVIDIA Cuda ON
\end{itemize}

In front-kick reference expansion runs, the code processed frames at approximately 10--16 frames per second in many segments, and the logs also contain repeated \textit{no detections} frames \cite{projectFrontKickLog,projectFrontKickStrictLog}. The seed batch run log similarly illustrates the non-negligible cost of continuous pose inference \cite{projectSeedLog}. This is consistent with the literature, which argues that real-time combat analysis is achievable only with careful management of resources and data quality \cite{quinn2024automation,quinn2025yolo}.

The implementation mitigates this through several engineering controls. The scripts expose configurable frame skipping, CPU thread limits, cooldown periods between batch jobs, preflight source checks, headless operation, and a \texttt{fast-mode} preset that disables expensive visual components \cite{projectActionRecognition,projectRunBatch}. Furthermore, CUDA enabled massive parallelism, allowing tasks to run much faster than on traditional CPUs. This has accelerated video processing significantly. 
These measures are not incidental; they are part of the project's answer to practical deployment constraints.

\section{Experimental Setup}
The uploaded Git Repository supports two related evaluation modes. The first is \textit{reference collection and calibration}, where canonical technique examples are captured and validated. The second is \textit{user evaluation}, where an athlete performs a target technique and the system compares that pose sequence against the calibrated reference bank.

Reference examples are stored as NumPy arrays of keypoint sequences in COCO-17 format \cite{projectActionRecognition}. The project documentation specifies a multi-angle folder structure for techniques, while the progress document shows that the intended capture plan consists of 52 total targets. The batch collection script aims by default for four examples per angle, subject to sufficient source diversity \cite{projectRunBatch}. This means the intended experiment design is systematic rather than ad hoc.

For evaluation, \texttt{calibrate\_label\_thresholds.py} requires validation examples for each technique split into \texttt{correct} and \texttt{incorrect} folders \cite{projectCalibrate}. Every validation clip is scored using the same best-reference matching routine used by the trainer, and the threshold is selected automatically from the observed score distribution. This is important because it means the acceptance rule is derived from validation evidence rather than hand tuning.

The present thesis can therefore extend the implemented system to the supervisor's requested comparison in a natural way. For an \textit{amateur versus professional} experiment, both groups can perform the same target techniques under comparable recording conditions, after which the system score, threshold decision, and score variance can be compared across groups. For an \textit{expert versus system} experiment, a coach or professional practitioner can provide an oral or manual score on a 1--10 scale for each attempt, and that score can then be compared against a normalized version of the system's internal score. The code supports the core metrics needed for these experiments.

\section{Results and Evaluation}

This section reports quantitative outputs generated by the implemented pipeline from four recorded runs stored under the structured storage directory (\texttt{data/runs}). The processed input videos were \texttt{WrongJab.mp4}, \texttt{RightJab.mp4}, \texttt{Jab\_Pro.mp4}, and \texttt{Jab\_Ama.mp4}. For all reported runs, the trainer threshold was 70/100 and metrics were aggregated from per-window scoring logs.

\subsection{Implemented Evaluation Capability}
The implemented system produces four useful outputs. First, it identifies the best matching reference angle for a given technique attempt \cite{projectActionRecognition}. Second, it computes a continuous technique score. Third, it converts that score into a binary correct-or-incorrect judgement using a score threshold (default or per-label calibrated, depending on configuration) \cite{projectActionRecognition,projectCalibrate}. Fourth, when an attempt is judged incorrect, it can generate a ghost-pose overlay to highlight the mismatch between the athlete and the reference \cite{projectActionRecognition}. Together, these outputs make the system suitable for coaching and formative assessment rather than only raw classification.

\subsection{Evidence from Logs}
The runtime logs provide direct evidence that the system was being used to capture and process references \cite{projectFrontKickLog,projectSeedLog}. In the expansion log, the system is shown as being armed, with the current reference techniques loaded and the video classifier disabled. This matches the intended workflow in the code: during reference capture, clean pose extraction and pose-based selection are more important than generic video classification.

The same logs also show intermittent sequences of frames with no person detections . This is a meaningful empirical finding. It supports the claim that robust combat-sports analysis depends on careful source selection, reference curation, and threshold calibration, because the raw detector output alone is not always stable enough across all frames of video.

In addition to runtime logs, the structured run outputs provide quantitative evidence. Across both runs, the system stored per-window scores, component metrics (cosine similarity, DTW distance, angle error, mean pose distance), selected reference angle, and generated feedback strings. This allows direct statistical reporting rather than anecdotal observations only.

\begin{table}[H]
\centering
\resizebox{\textwidth}{!}{
\begin{tabular}{lcccccc}
\toprule
Source clip & Windows & Accepted & Acceptance rate (\%) & Program Score & Trainer Score & Std. dev. \\
\midrule
\texttt{Jab\_Pro.mp4} & 564 & 465 & 82.4 & 91.81 & 85 & 10.41  \\
\texttt{Jab\_Ama.mp4} & 632 & 354 & 56.0 & 62.07 & 60 & 21.18  \\
\texttt{RightJab.mp4} & 316 & 217 & 75.7 & 83.03 & 80 & 9.87  \\
\texttt{WrongJab.mp4} & 296 & 132 & 44.6 & 42.11 & 60 & 18.27  \\
\bottomrule
\end{tabular}
}
\caption{Run-level trainer acceptance summary from structured metrics logs.}
\label{tab:run_level_summary}
\end{table}

The \texttt{Jab\_Pro.mp4} run produced both the highest mean score and the highest acceptance rate among the four runs, while \texttt{WrongJab.mp4} produced the lowest mean score. The \texttt{Jab\_Ama.mp4} run showed the highest score variance, indicating unstable execution quality across windows. These patterns are consistent with the intended behavior of the scoring pipeline under a fixed threshold.

\subsection{Clip-Level and Trainer-Only Professional/Amateur Comparison}
The four images illustrate a qualitative comparison between a stronger and a weaker jab execution. At this stage, the quantitative evidence is clip-level rather than a full participant study. The system behavior is consistent with the intended trainer logic: stronger attempts tend to remain above threshold with fewer corrective prompts, while weaker attempts fall below threshold and trigger ghost-pose guidance.

At frame level, this behavior can be seen directly in the stored metrics: windows above 70/100 are marked correct, while lower-score windows activate corrective feedback such as \textit{Extend your punching arm more} and \textit{Keep your guard hand higher}. This aligns with the visible differences in guard position and extension quality in Figure~\ref{fig:jab_comparison}.

\begin{figure}[h!]
\centering

% Top row
\begin{minipage}{0.48\textwidth}
    \centering
    \includegraphics[width=\linewidth, height=6cm, keepaspectratio]{Jab_FP_pro.png}
\end{minipage}
\hfill
\begin{minipage}{0.48\textwidth}
    \centering
    \includegraphics[width=\linewidth, height=6cm, keepaspectratio]{Jab_FP_ext_Pro.png}
\end{minipage}

\vspace{0.5cm}

% Bottom row
\begin{minipage}{0.48\textwidth}
    \centering
    \includegraphics[width=\linewidth, height=6cm, keepaspectratio]{Jab_Ama_Guard_Good.png}
\end{minipage}
\hfill
\begin{minipage}{0.48\textwidth}
    \centering
    \includegraphics[width=\linewidth, height=6cm, keepaspectratio]{Jab_Ama_Guard.png}
\end{minipage}

\caption{Comparison of jab techniques. The top row shows professional execution and extension, while the bottom row illustrates guard positioning in correct and incorrect forms.}
\label{fig:jab_comparison}

\end{figure}


\begin{table}[H]
\centering
\begin{tabular}{lccc}
\toprule
Clip & Mean Score & Pass Rate & Score Std. Dev. \\
\midrule
Clip A (\texttt{WrongJab\_Input.mp4}) & 69.00 & 54.84\% & 8.27 \\
Clip B (\texttt{Jab.mp4}) & 60.03 & 22.44\% & 9.87 \\
\bottomrule
\end{tabular}
\caption{Comparison of two recorded user clips using the implemented trainer pipeline.}
\label{tab:amateur_pro_template}
\end{table}

Under this sample, Clip A achieved both higher mean score and lower variance than Clip B, suggesting more stable execution under the current references and threshold. A full professional-vs-amateur comparison should aggregate multiple attempts per participant and then compare group-level means and variances.

For the trainer-only professional/amateur comparison requested by the project supervisor, the program score itself is used as the evaluation variable. The conversion to a 1--10 scale is:
\[
\text{Program}_{1\text{--}10} = \frac{\text{Program}_{0\text{--}100}}{10}.
\]

Table~\ref{tab:trainer_only_pro_ama} reports trainer-only summary statistics for the dedicated professional and amateur jab runs.

\begin{table}[H]
\centering
\begin{tabular}{lccc}
\toprule
Run & Program Mean (0--100) & Program Mean (1--10) & Acceptance Rate (\%) \\
\midrule
Professional (\texttt{Jab\_Pro\_Done.mp4}) & 60.81 & 6.08 & 26.42 \\
Amateur (\texttt{Jab\_Ama\_Done.mp4}) & 62.07 & 6.21 & 37.02 \\
\bottomrule
\end{tabular}
\caption{Trainer-only professional versus amateur comparison using program scores.}
\label{tab:trainer_only_pro_ama}
\end{table}

In this sample, the trainer-only metrics do not yet show a clean professional-higher/amateur-lower separation. This should be interpreted as an indicator that additional controlled recordings and per-technique threshold calibration are still required before drawing proficiency conclusions.

\subsection{Threshold Calibration Results}
The project includes threshold calibration tooling \cite{projectCalibrate}; however, no exported per-technique calibration file was present in the current repository snapshot. For the reported runs, the trainer operated with the configured default threshold of 70.0 for \texttt{jab}.

Therefore, the present quantitative results should be interpreted as fixed-threshold evaluation, not final per-technique calibrated performance. Table~\ref{tab:threshold_template} reflects the currently available evidence.

\begin{table}[H]
\centering
\begin{tabular}{lcccccc}
\toprule
Technique & Threshold & F1 & Bal. Acc. & Precision & Recall & Samples \\
\midrule
Jab (\texttt{WrongJab\_Input.mp4}) & 70.0 & -- & -- & -- & -- & 155 \\
Jab (\texttt{Jab.mp4}) & 70.0 & -- & -- & -- & -- & 967 \\
Jab (\texttt{Jab\_Pro\_Done.mp4}) & 70.0 & -- & -- & -- & -- & 564 \\
Jab (\texttt{Jab\_Ama\_Done.mp4}) & 70.0 & -- & -- & -- & -- & 1364 \\
\bottomrule
\end{tabular}
\caption{Available threshold evidence from current runs. Full calibration metrics require running the calibration pipeline on labeled validation sets.}
\label{tab:threshold_template}
\end{table}

\subsection{Component Metric Analysis}
To understand score behavior, component means from run summaries were compared. For \texttt{WrongJab\_Input.mp4}, DTW and mean pose distance were both approximately 0.43. For \texttt{Jab.mp4}, they increased to approximately 2.15 and 2.19. For \texttt{Jab\_Pro\_Done.mp4}, they were approximately 1.67 and 1.70. For \texttt{Jab\_Ama\_Done.mp4}, they were approximately 2.28 and 2.39. Because the final score combines cosine, DTW, angle, and pose-distance terms, higher temporal and geometric distance generally lowers acceptance.

Most runs matched \texttt{left45\_02} most frequently, while \texttt{Jab\_Ama\_Done.mp4} most often matched \texttt{jab\_left145\_01}. Across all runs, the most common corrective cues were \textit{Extend your punching arm more} and \textit{Keep your guard hand higher}, consistent with recurring extension and guard-maintenance errors.

\section{Discussion}
The implemented system shows that a pose-based reference-matching approach is a viable design for combat sports technique evaluation. This is an important distinction from more conventional HAR systems. Instead of only recognizing what action occurred, the system attempts to determine how well the action was performed. That shift makes the project more useful for coaching, but it also introduces stricter requirements on data quality, pose stability, and threshold calibration.

The finalized trainer-only results also show an important practical point: in the current configuration, the professional and amateur sample runs do not yet produce a stable skill-ordering separation. This should not be interpreted as failure of the approach, but as evidence that score behavior is still sensitive to recording setup, reference selection, and threshold choice. In other words, the current results validate pipeline functionality, while also showing that stronger experimental controls are still needed before making strong proficiency claims.

One of the main strengths of the project is its modularity. Reference collection, candidate filtering, calibration, scoring, visualisation, and batch processing are separated into dedicated scripts, which makes the workflow transparent and extensible \cite{projectScrape,projectFilter,projectGenerate,projectRunBatch,projectCalibrate}. Another strength is the explicit accommodation of camera angle through multi-angle references. This is a practical and well-justified answer to a genuine source of error in martial arts analysis.

The main limitations are equally clear. First, the system depends heavily on the quality of the reference bank. Weak or poorly segmented references will produce misleading scores. Second, the detector still produces occasional gaps in person detection, especially in uncontrolled web footage \cite{projectFrontKickLog,projectFrontKickStrictLog}. Third, the final quality of the trainer depends on having enough validation data per technique to calibrate thresholds meaningfully. When such data is missing, the system must fall back to default thresholds, which is a practical compromise but not an ideal research endpoint \cite{projectCalibrate}.

\section{Conclusion and Future Work}
This thesis presented a computer-vision-based system for combat sports technique analysis built around pose estimation, temporal pose sequence comparison, multi-angle references, and calibrated scoring thresholds. The implemented project goes beyond basic action classification by evaluating the quality of technique execution and providing feedback through thresholds and ghost-pose overlays. This makes it relevant not only as a recognition system but also as a training support tool.

The development process revealed several major challenges: collecting reliable reference data, handling viewpoint variation, segmenting short combat actions correctly, transforming raw similarity into meaningful correctness decisions, and maintaining acceptable runtime behaviour on practical hardware. Each of these challenges led to concrete design choices in the codebase, including semi-automated reference collection, multi-angle reference storage, event-centred sequence extraction, per-technique threshold calibration, and runtime controls.

Future work should focus on three areas. First, the reference library should be expanded so that all planned techniques and angles are represented by multiple high-quality examples collected under controlled capture conditions. Second, per-technique threshold calibration should be completed on labeled validation sets so acceptance decisions are data-driven rather than default-threshold driven. Third, once protocol control is tightened, an expert-versus-program study should be executed with synchronized 1--10 labels to quantify human-system agreement on the same attempts.

\printbibliography

\end{document}